\section{Engineering items}\label{sec:engineering}

Each item retains a motivation from the original \code{33cec8b} planning
snapshot, a revised design, and prospective exit evidence. These descriptions
are not a current-checkout defect inventory. The integrated semantic boundary
in Section~\ref{sec:semantic-boundaries} and the phased plan in
Section~\ref{sec:roadmap} qualify the original implementation suggestions.

\subsection{E1. Results as compilable definitions}

\paragraph{Historical motivation.} The profiled REPL bound candidates as
\code{it1}, \code{it2}, \dots{} with \code{addAndCompile}. A candidate built
from a recursor, such as the list sum
\begin{lstlisting}
fun a => List.rec Nat.zero (fun head tail tail_ih => head.add tail_ih) a
\end{lstlisting}
was accepted by the kernel and then rejected by the code generator (``does not
support recursor \code{List.rec}''). That run displayed a correct answer
followed by errors and did not provide an evaluable published result.

\paragraph{Design.} Separate the certified term from the published definition.
\begin{enumerate}
\item Preserve the exact certified program and original contract proof. Export
them into a staging environment, abstracting the original local telescope and
universe parameters, and recheck the resulting declarations. Separating
\code{addDecl} from compilation avoids conflating two operations; it does not
prove that declaration transport or binding cannot fail.
\item Construct an executable presentation through Lean's equation compiler.
Where a recursor plan supports a pattern-matching definition, elaborate that
definition, compile it, and prove its correspondence with the logical term
or prove the original contract directly for the executable definition.
Try definitional equality first; use induction or the checked equation
theorems when required. A source-level \code{match} rendering by itself is not
a compilation or correspondence proof.
\item Publish the checked executable definition when all promised obligations
hold. An \code{@[implemented_by]} route requires its own audited backend and
correspondence assumptions; the attribute is not the theorem connecting two
programs. If only the logical artifact is available, label export and
compilation statuses explicitly, and offer a \code{noncomputable} presentation
only if that declaration has itself checked. It does not satisfy an executable
publication request.
\end{enumerate}
Constant-motive non-indexed folds are an appropriate first supported fragment;
arbitrary recursor syntax can still contain nested eliminations or scope and
sharing issues. Nonconstant motives and dependent arguments require the
realization obligations in Sections~\ref{sec:dependent} and~\ref{sec:recursion}.
Use version-pinned native elaboration interfaces and retain explicit unsupported
outcomes rather than asserting a universal mechanical translation.

\paragraph{Closing test.} Probe 6 exports and compiles without error and
\texttt{\#eval it1 [1,2,3]} prints 6. Every supported recursive candidate has
separate logical, exported and compiled receipts and a readable source
presentation whose elaboration is rechecked. Printing a \code{match} is a
presentation test, not a substitute for those gates.

\subsection{E2. Frontends}

The original plan began with the REPL and \texttt{\#leant2}. A shared query
interface can support additional frontends, but dependent local contexts,
universe preservation, hole ownership and replay make their adapters substantive
correctness boundaries rather than merely syntax wrappers.

\begin{description}
\item[Term elaborator] \texttt{synth\%} elaborates to a certified candidate
under the declared ranking. Postpone unresolved expected-type dependencies
when necessary; do not specialize the original universe parameters or assign
foreign holes merely to make the query easier. Timeout without a certified
candidate remains an incomplete result.
\item[Tactic] \code{leant2} closes a goal whose target is data or a proposition, reporting the term through \code{Try this}. It competes with \code{exact?}; its advantage is construction (case analysis, recursion, contracts) rather than lookup.
\item[Sketches] A body such as
\code{def f xs := List.foldr ?step ?init xs}, together with an explicitly
associated contract, supplies a partial program and carrier hint. Distinguish
owned holes from pre-existing elaborator metavariables, close each typed local
interface, preserve dependencies and apply the final original contract to the
whole completed program. Sketch surface syntax and hole elaboration must be
specified and tested, rather than inferred from the existence of an internal
metavariable search state.
\item[Code action] Offer ``synthesize body'' at a \code{sorry} with explicitly
associated tests. Recognize and check each test's relation to the declaration;
textual proximity alone does not identify the intended contract. Replay the
replacement source in its original declaration context.
\item[Library API] A documented \code{Leant2.synthesize : Query -> MetaM Outcome} with the result algebra, for use from other metaprograms, including Forge.
\end{description}

\paragraph{Closing test.} Each frontend has native success, scope, policy and
interruption cases. The prospective sketch target is to solve all thirteen
stretch tasks with their reference carriers and unchanged contracts. Record
body-search failures separately: supplying a carrier does not establish that
transition and finishing synthesis will succeed.

\subsection{E3. Receipts, replay and resumable search}

The unified proposal specifies candidate receipts containing query, environment,
program, contract proof, axioms and ledger. Preserve proof replay separately
from logical search replay and executable publication. A resumable incomplete
outcome also needs the original problem, pending alternatives and proofs,
semantic state, expansion cursor, policy/model epoch and monotone work charges.
Section~\ref{sec:search-deadlines} gives the full checkpoint contract.

A deterministic lane/depth/alternative path is one possible reconstruction
format if all effects, ordering and external proposals needed to reproduce it
are pinned. It is not sufficient by itself for branch caches, delayed
constraints or dynamic priorities. Compare reconstruction cost and native
persistent snapshots experimentally; neither is universally cheaper or more
robust by construction.

\paragraph{Closing test.} Cutting and resuming at recorded logical checkpoints
reproduces the candidate and event sequence of uninterrupted execution to the
same final logical prefix. Separately measure reconstruction and startup time.
Three two-second executions need not reach the same prefix as one six-second
execution, as Proposition~\ref{prop:search-deadline} explains.

\subsection{E4. Substrate cost}

The historical profile left thirty percent of contract-query time unattributed
(Table~\ref{tab:profile}). The plan is measurement first.
\begin{enumerate}
\item Build the executable with a sampling profiler attached (on Windows, ETW; on Linux, \code{perf}) and read the C call stacks of the Lean runtime. This distinguishes allocation, reference counting, persistent-map operations and \code{instantiateMVars} traversals without guessing.
\item Re-profile the repeated scans and instantiations identified by the
historical source inspection. Cache per-local facts only with their complete
assignment, universe, transparency and instance dependencies; introducing a
local once does not make its type immutable while metavariables remain open.
Keep branch-valid caches in the goal/state representation, with invalidation
and rollback tests from Section~\ref{sec:evaluation}.
\item Replace linear scans by indexes (item A5).
\end{enumerate}
\paragraph{Closing test.} The original targets of unattributed time below
10\% and node cost below 1 ms on the Church probes remain measurable
hypotheses. Report query mix, inclusive/exclusive accounting, cache overhead
and retained memory; moving time between counters is not a speedup.

\subsection{E5. Parallel lanes}

Whole-lane parallelism is a candidate optimization when each worker owns its
semantic state and the shared environment is frozen. Audit all external
references, environment mutations and services before declaring the workers
independent; an \code{IO.Ref} does not establish safe ownership. Shared work
accounting and cancellation require an explicit synchronization protocol.

For deterministic logical replay, give jobs stable identities, fixed work
quanta and a canonical commit order. Workers may speculate, but admit completed
epochs by that order. First-found cancellation is a separate opportunistic
policy: ranking its survivors afterwards cannot recover candidates lost to
the race. Preserve its actual event trace without claiming identical timed
answer sets. The detailed boundary is in Section~\ref{sec:search-deadlines}.

OR-parallel alternatives and carrier proposals follow the same ownership
discipline. AND-parallel subgoals additionally need the factorization
conditions of Section~\ref{sec:search-factorization}. Measure memory, duplicated
work, publication latency and lost useful work on cancellation, not only CPU
utilization.

\paragraph{Closing test.} Compare sequential and parallel lanes at matched
logical work and separately under a common interactive budget, preserving the
same certified outcomes and replaying completed checkpoints. Faster classical
or refutation results are an empirical target, not an automatic consequence
of scheduling tasks concurrently.

\subsection{E6. A decision procedure for propositional targets}

For the reified implicational-conjunctive-disjunctive propositional fragment,
intuitionistic provability is decidable and contraction-free LJT calculi give
terminating search~\cite{dyckhoff92}. A checked finite Kripke countermodel
establishes nonvalidity in that calculus. It does not by itself prove a Lean
negation for an arbitrary contextual target whose atoms have an interpretation.
Positive reconstruction and any negative transfer both need the reification
bridge and complete treatment of the relevant assumptions specified in
Section~\ref{sec:semantic-boundaries}. Likewise, a SAT result about the reified
classical formula needs its appropriate checked certificate and interpretation.

Such results can justify scheduling constructive and policy-permitted
classical work for the supported fragment. They do not decide arbitrary Lean
propositions merely because \code{Prop} occurs in their type. General
intuitionistic propositional validity is PSPACE-complete~\cite{statman79};
small practical instances and latency must be measured.

\paragraph{Closing test.} Inventory which \code{synth-basic} and
\code{synth-manual} queries satisfy the supported reification assumptions.
Each covered query receives a checked positive reconstruction or a precisely
scoped negative certificate; unsupported contextual queries remain explicit.
The original under-50-ms target is prospective and reported separately from
certificate validity.

\subsection{E7. A wider proof portfolio}

The historical portfolio \code{first | rfl | decide | simp | omega} motivated
a transactional adapter with typed replies and per-tier budgets. Evaluate
available \code{grind}, \code{try?}, arithmetic simplification and bounded
instantiation adapters at the pinned toolchain; names alone do not establish
integration or cost. Local-context proof obligations and induction hypotheses
need the richer interface of Section~\ref{sec:contracts}.

For $\forall x:\tau,\ P\,x$, a concrete well-typed input with a checked
proof of $\neg P(x)$ refutes the candidate. Enumerating a few values can find
such a counterexample without a small-model theorem; passing those tests does
not prove the universal contract. Finite exhaustive validation requires an
enumeration coverage proof and exact evaluation. Retain unknown evaluation or
proof outcomes and the axiom policy for negative evidence. Schedule testing
and proving fairly rather than making an arbitrary finite test prefix the
sole gate to all proof work.

\paragraph{Closing test.} A suite of twenty quantified contracts over lists and naturals (length preservation, idempotence, involution) where wrong candidates are refuted by instantiation in under 5 ms each.

\subsection{E8. Benchmarks beyond Leant}

The historical scored Leant corpus was saturated; that does not establish
coverage of its unscored stretch tasks or of broader synthesis. Port suitably
scoped tasks from Myth~\cite{myth}, $\lambda$\textsuperscript{2}~\cite{lambda2},
Smyth~\cite{smyth}, Burst~\cite{burst}, Trio~\cite{trio}, and the type-only
Agda/Coq studies~\cite{agsy,sauto}, deduplicating semantic families and recording
each source's actual grammar, supplied helpers, examples and recursion scope.
Add Lean tasks by erasing small core \code{List}, \code{Option} and \code{Nat}
definitions and selecting suitable lemmas as explicit contracts. A published
score in another language is not a matched Leant2 baseline by default.
The historical stretch cases and failing probes provide initial tasks, with
separate carrier/motive-given diagnostics. Use the complete timing, evidence
and ablation matrix in Section~\ref{sec:roadmap}; report certification and
executable publication separately from first candidate discovery.
